%!TEX root = ../../main.tex
%% ===============================================================
%% 10.3 코드--논문 감사
%% 파일: chapters/10_discussion/03_audit.tex
%% 상위: chapters/10_discussion/chapter.tex
%% 정의 라벨: sec:disc-audit
%% 참조 라벨: tab:audit, tab:baseline, tab:conditional
%% 포함 객체: tables/tab_audit
%% 인용: -
%% 상태: 초안 (docs/OUTLINE.md 참고)
%% ===============================================================

\section{코드--논문 감사}
\label{sec:disc-audit}

학술 논문의 방법 기술과 실제 코드 사이의 불일치는 재현성의 주요 위협이다. 본 연구는 최종본 제출 전에 코드와 논문을 항목별로 대조하는 감사를 수행하였고, 발견된 모든 불일치를 논문의 기술에 맞게 코드를 고치는 방향으로 해소하였다. 표 \ref{tab:audit}에 그 내역을 요약한다.

%% ---------------------------------------------------------------
%% table 객체: tab:audit
%% 파일: tables/tab_audit.tex
%% 사용처: chapters/10_discussion/03_audit.tex  (%% ---------------------------------------------------------------
%% table 객체: tab:audit
%% 파일: tables/tab_audit.tex
%% 사용처: chapters/10_discussion/03_audit.tex  (%% ---------------------------------------------------------------
%% table 객체: tab:audit
%% 파일: tables/tab_audit.tex
%% 사용처: chapters/10_discussion/03_audit.tex  (\input{tables/tab_audit})
%% 규칙: 캡션·라벨·내용을 이 파일 안에서만 관리한다. 수치는 config/numbers.tex 매크로를 쓴다.
%% ---------------------------------------------------------------
\begin{table}[htbp]
  \centering
  \caption{사전 감사에서 수정된 코드--논문 불일치}
  \label{tab:audit}
  \small
  \begin{tabularx}{\linewidth}{@{}lXX@{}}
    \toprule
    항목 & 발견 & 수정 \\
    \midrule
    A3 지름 분할 & 단일 연결(4{,}000 단위)로 분할하여 사슬형 군집의 지름이 상한을 초과할 수 있었음 & 완전 연결·지름 상한 분할로 교체 \\
    A1 온셋 검증 & ``생존 2 명(지도 전체)''과 ``반경 내 2 명(생존 무관)''을 별개 조건으로 검사 & 생존 $\wedge$ 반경 내를 단일 논리곱으로 검사 \\
    A2 인접 병합 & 15 초 / 2{,}000 단위 병합이 설정만 있고 실행되지 않음 & 병합 단계 구현 \\
    B1 특수 킬 & \texttt{CHAMPION\_SPECIAL\_KILL}을 두 번째 킬로 계수 & 보너스만 반영 \\
    MLP 경로 & 헤드라인 러너의 MLP가 95 차원 거시 시퀀스를 사용 & 2{,}980 차원 정합 입력으로 재경로 \\
    \bottomrule
  \end{tabularx}
\end{table}
)
%% 규칙: 캡션·라벨·내용을 이 파일 안에서만 관리한다. 수치는 config/numbers.tex 매크로를 쓴다.
%% ---------------------------------------------------------------
\begin{table}[htbp]
  \centering
  \caption{사전 감사에서 수정된 코드--논문 불일치}
  \label{tab:audit}
  \small
  \begin{tabularx}{\linewidth}{@{}lXX@{}}
    \toprule
    항목 & 발견 & 수정 \\
    \midrule
    A3 지름 분할 & 단일 연결(4{,}000 단위)로 분할하여 사슬형 군집의 지름이 상한을 초과할 수 있었음 & 완전 연결·지름 상한 분할로 교체 \\
    A1 온셋 검증 & ``생존 2 명(지도 전체)''과 ``반경 내 2 명(생존 무관)''을 별개 조건으로 검사 & 생존 $\wedge$ 반경 내를 단일 논리곱으로 검사 \\
    A2 인접 병합 & 15 초 / 2{,}000 단위 병합이 설정만 있고 실행되지 않음 & 병합 단계 구현 \\
    B1 특수 킬 & \texttt{CHAMPION\_SPECIAL\_KILL}을 두 번째 킬로 계수 & 보너스만 반영 \\
    MLP 경로 & 헤드라인 러너의 MLP가 95 차원 거시 시퀀스를 사용 & 2{,}980 차원 정합 입력으로 재경로 \\
    \bottomrule
  \end{tabularx}
\end{table}
)
%% 규칙: 캡션·라벨·내용을 이 파일 안에서만 관리한다. 수치는 config/numbers.tex 매크로를 쓴다.
%% ---------------------------------------------------------------
\begin{table}[htbp]
  \centering
  \caption{사전 감사에서 수정된 코드--논문 불일치}
  \label{tab:audit}
  \small
  \begin{tabularx}{\linewidth}{@{}lXX@{}}
    \toprule
    항목 & 발견 & 수정 \\
    \midrule
    A3 지름 분할 & 단일 연결(4{,}000 단위)로 분할하여 사슬형 군집의 지름이 상한을 초과할 수 있었음 & 완전 연결·지름 상한 분할로 교체 \\
    A1 온셋 검증 & ``생존 2 명(지도 전체)''과 ``반경 내 2 명(생존 무관)''을 별개 조건으로 검사 & 생존 $\wedge$ 반경 내를 단일 논리곱으로 검사 \\
    A2 인접 병합 & 15 초 / 2{,}000 단위 병합이 설정만 있고 실행되지 않음 & 병합 단계 구현 \\
    B1 특수 킬 & \texttt{CHAMPION\_SPECIAL\_KILL}을 두 번째 킬로 계수 & 보너스만 반영 \\
    MLP 경로 & 헤드라인 러너의 MLP가 95 차원 거시 시퀀스를 사용 & 2{,}980 차원 정합 입력으로 재경로 \\
    \bottomrule
  \end{tabularx}
\end{table}


수정의 순효과는 말뭉치가 \numEngagementsPaper{} 건에서 \numEngagementsCorrected{} 건으로 \corpusReductionPct{}\% 줄어든 것(일방적 정리와 중복 재교전이 제거됨)과 LightGBM 시험 AUC가 $\aucLGBM{}$에서 약 $\aucLGBMCorrected{}$로 $\aucDropTotal{}$ 낮아진 것이다(국소화 $\aucDropLoc{}$, 라벨 $\aucDropLabel{}$). 계열 간 순서(테이블 $\gg$ 신경망 표현)와 조건부 예측 가능성의 결론은 바뀌지 않았다. 본 논문은 심사를 거친 수치를 본문에 유지하고 이 절에서 수정 후 수치를 명시한다. \todo{학위논문 최종본에서는 수정된 코드로 전체 실험(3 시드 $\times$ 2 단계)을 재실행하여 표 \ref{tab:baseline}--\ref{tab:conditional}을 갱신하는 것을 권장한다.}
